\chapter{The Platform Map}\label{ch:pl:the-platform-map}

At 07:40 the overnight risk report, the P\&L and a researcher's backtest disagree about one position. Three teams spend the
morning finding out why. The vendor's corporate-action file, which normally lands at 19:30, came at 21:30 the night before; the
risk batch, which starts at 21:00 by the clock, had already read the old security master and the positions built on it; the P\&L
job, which starts when its inputs are ready, read the new ones twenty minutes later. Every system did what it was built to do.
What was missing was a picture, kept as data and checked by a program, of which system reads what, from whom, and when. This
chapter draws that picture for the miniature firm built across the series, and turns it into a tool (\texttt{firm.platmap})
that answers the questions the three teams spent their morning on: what a late file reaches, which job reads an input before it
is complete, and how late each input may arrive before the morning's reports miss their deadline.

\section{The systems of a trading firm}

A trading firm is a chain of systems that turn market data and decisions into positions, and positions into money that has
been counted, checked and reported. Books~1 to 14 of this series built most of the links one at a time: the exchange simulator
and its feed (Book~10), the feed handler, order gateway and risk gate (Book~13), the backtesters and the research workflow
(Book~7), the pricing library (Book~5), the risk engine (Book~6), the machine-learning platform (Book~12). This book is about
the platforms that connect them, and it starts with the map.

\begin{definition}[Front-to-back flow]\label{def:pl:the-platform-map:ftb}
The \emph{front-to-back flow}\index{front-to-back flow} of a trading firm is the path that a trade and the data about it take
through the firm's systems: market and reference data; the decision and the order; the execution; the position and its P\&L;
the risk it adds; its booking, confirmation, settlement and reconciliation; and the controls and reports that close the day.
\end{definition}

The industry names the stretches of that path after the teams that historically ran them: the front, middle and back office,
which One Quant Book~16 (chapter~13) defines as organisational units. This book cares about the systems, whoever runs them.
\Cref{fig:pl:the-platform-map:layers} shows the miniature firm's eighteen systems in eight layers, and the direction in which
the day's data move through them.

\begin{omfigure}
\begin{tikzpicture}[font=\scriptsize,>=Stealth,
  sys/.style={draw,rounded corners=2pt,fill=omDef!10,minimum height=0.55cm,inner sep=2pt,align=center},
  ext/.style={sys,fill=gray!12},
  lay/.style={font=\scriptsize\itshape,anchor=east,text=gray!70!black}]
\foreach \y/\name in {0/external,-1/data,-2/trading,-3/positions,-4/risk,-5/post-trade,-6/control}
  {\node[lay] at (-0.2,\y) {\name};}
\node[ext] (ex) at (1.5,0) {exchanges};
\node[ext] (dv) at (4.4,0) {data vendor};
\node[ext] (cb) at (8.8,0) {clearing broker};
\node[sys] (tc) at (1.0,-1) {tick capture};
\node[sys] (ts) at (2.9,-1) {tick store};
\node[sys] (rd) at (5.0,-1) {reference data};
\node[sys] (ss) at (7.3,-1) {snapshot service};
\node[sys] (og) at (1.9,-2) {strategies and gateway};
\node[sys] (ps) at (1.9,-3) {position service};
\node[sys] (pn) at (4.2,-3) {P\&L service};
\node[sys] (pl) at (4.2,-4) {pricing library};
\node[sys] (rg) at (6.8,-4) {risk grid};
\node[sys] (tcap) at (1.9,-5) {trade capture};
\node[sys] (rc) at (8.8,-5) {reconciliation};
\node[sys] (pc) at (1.9,-6) {product control};
\node[sys] (sv) at (3.9,-6) {surveillance};
\node[sys] (rr) at (6.8,-6) {regulatory reporting};
\node[sys,fill=omMeth!12,minimum height=2.4cm,text width=1.5cm] (rp) at (10.9,-2.7) {research platform};
\draw[->] (ex) -- (tc);
\draw[->] (tc) -- (ts);
\draw[->] (dv) -- (rd);
\draw[->] (og) -- (ps);
\draw[->] (ps) -- (pn);
\draw[->] (pl) -- (pn);
\draw[->] (pl) -- (rg);
\draw[->] (ps) -- (tcap);
\draw[->] (cb) -- (rc);
\draw[->] (ts.south) |- (rp.north west |- 0,-1.55);
\draw (rd.south) -- (rd.south |- 0,-1.55);
\draw[->] (ss.south) -- (ss.south |- rg.north);
\end{tikzpicture}

\omcaption{The miniature firm's systems by layer, from the venues at the top to the controls at the bottom; the research platform
reads the tick store and the reference data beside the trading path. Arrows show the main direction of the day's data; the full
set of flows is the map's data (\texttt{pl\_platmap}), not the drawing.}\label{fig:pl:the-platform-map:layers}
\end{omfigure}

Read from the top, the layers are the order in which a trade's data are born. Market data and reference data exist before any
decision; the order gateway sends orders that the venue fills; the position service keeps what the fills add up to; the P\&L
service values it with the pricing library's marks; the risk grid revalues it under scenarios; trade capture books it and
post-trade systems confirm, settle and reconcile it; product control, surveillance and regulatory reporting check and report
all of the above. A research platform sits to the side: it reads the same tick store and reference data to build the next
strategy, and its outputs come back as parameters and models that the trading layer loads (chapter~16).

\begin{example}[The miniature firm in numbers]\label{ex:pl:the-platform-map:numbers}
The map of the miniature firm (\texttt{pl\_platmap}) has 18 systems, 22 datasets and 15 end-of-day jobs. Seven datasets
arrive from outside: the day's fills, the day's market-data capture, the official closing prices, a market-data snapshot of
curves and volatilities, the vendor's reference-data and corporate-action files, and the clearing broker's statement. Every
other dataset is produced by exactly one job from others.
\end{example}

\section{Data flows and their clocks}

The difficulty of a platform is not the number of systems but the number of clocks. Market data and fills arrive continuously
during the session; the official closing prices come some minutes after the close; vendor files come in the evening at times
the vendor does not guarantee; the clearing broker's statement comes when the clearing house has finished its own day. Each job
of the end-of-day chain runs on one of two kinds of clock: it starts when its inputs are ready, or it starts at a time on the
wall.

\begin{definition}[Data-flow map]\label{def:pl:the-platform-map:map}
A \emph{data-flow map}\index{data-flow map} is a description, kept as data, of a firm's systems, the datasets each one holds,
and the jobs that read datasets and write others, each job with its duration and its clock (dependency-driven or scheduled at a
fixed time), each dataset with its \omterm{def:pl:the-platform-map:sor}{system of record} and, for external inputs, its arrival time.
\end{definition}

Because it is data, a map can be checked (every dataset has one owner, the jobs form no cycle), traversed (what does this file
reach?) and simulated (when will each report be ready tonight?). A diagram drawn once for a presentation does none of these, and
is out of date by the next release.

\begin{omfigure}
\begin{tikzpicture}[font=\scriptsize,>=Stealth]
\draw[->] (0,0) -- (12.5,0) node[right]{time};
\foreach \x/\t in {0/09:30,3.63/16:00,5.58/19:30,6.70/21:30,9.77/03:00,12.0/07:00}
  {\draw (\x,0.08) -- (\x,-0.08) node[below]{\t};}
\fill[omDef!25] (0,0.3) rectangle (3.63,0.75); \node at (1.8,0.525) {stream: ticks, fills};
\fill[omMeth!25] (3.67,0.3) rectangle (10.28,0.75); \node at (7.0,0.525) {end-of-day batch};
\fill[omThm!20] (6.93,0.95) rectangle (10.28,1.4); \node at (8.6,1.175) {overnight research};
\draw[omThm,thick] (12.0,0.2) -- (12.0,1.6) node[above]{reports due};
\draw[->,thick] (3.63,1.55) node[above]{close} -- (3.63,0.8);
\draw[->] (5.58,1.55) node[above]{vendor files} -- (5.58,0.8);
\draw[->] (6.70,2.15) node[above]{clearing statement} -- (6.70,0.8);
\end{tikzpicture}

\omcaption{The clocks of a trading day. Ticks and fills stream during the session; the \omterm{def:pl:the-platform-map:batch}{end-of-day batch} starts at the close and
waits for files whose arrival times it does not control; overnight research runs after the batch has produced its inputs; the
morning reports are due at 07:00. Times are the miniature firm's nominal ones.}\label{fig:pl:the-platform-map:clocks}
\end{omfigure}

\Cref{lst:pl:the-platform-map:jobs} is the end-of-day chain of the miniature firm as the map records it: each job with its
inputs, outputs and duration in minutes. The risk batch takes an optional fixed start; everything else starts on its inputs.

\omcode{code/platforms/01-the-platform-map/python/pl_platmap.py}{61}{81}{The end-of-day jobs of the miniature firm: inputs,
outputs, duration in minutes and, for the risk batch, an optional fixed start.}\label{lst:pl:the-platform-map:jobs}

\section{Systems of record and golden sources}

A number exists in several places at once: the closing price of a stock is in the exchange's feed, in the tick store, in the
pricing library's snapshot, in the P\&L service's marks and in every spreadsheet that copied it. When two of those copies
differ, the firm needs a rule that says which one is right.

\begin{definition}[System of record, golden source]\label{def:pl:the-platform-map:sor}
The \emph{system of record}\index{system of record} of a dataset is the one system whose copy is authoritative: corrections are
made there and nowhere else, and every other copy is derived from it. A \emph{golden source}\index{golden source} is a system of
record that the firm designates for a kind of data used across many systems (instruments, counterparties, closing marks), so
that every consumer reads from it or from a copy that the map shows to be derived from it.
\end{definition}

\begin{proposition}[One owner, one answer]\label{prop:pl:the-platform-map:oneowner}
If every dataset has exactly one \omterm{def:pl:the-platform-map:sor}{system of record}, every copy is derived through the map from that system, and two reports read
the same version of a dataset, then they agree about it.
\end{proposition}

\begin{proof}
Each report's value of the dataset is a function of the version it read, applied along the map's derivations; the same version
through the same derivations gives the same value. A disagreement therefore requires two systems of record (two versions that
are not derived from one another), a copy made outside the map (a derivation nobody declared), or two versions read at
different times.
\end{proof}

The proof is short because the proposition is almost a definition, and its value is in the three ways it can fail. The first
two are structural and a program can find them in the map; the third is about time, and the rest of this chapter is about it.
\Cref{lst:pl:the-platform-map:validate} is the structural check. On the map as first drawn it finds that two systems claim the
marks, the pricing library and product control.

\begin{example}[Two claims on the marks]\label{ex:pl:the-platform-map:marks}
The pricing library computes the end-of-day marks from the closing prices and the curve and volatility snapshot; product
control independently verifies them against broker quotes and consensus data (One Quant Book~6, chapter~27). Both call their
output ``the marks''. They are two datasets: the marks, owned by the pricing library, and the verified marks with their
adjustments, owned by product control. The map's fix is to name the second dataset and make its dependence on the first
explicit; the check then passes, and a disagreement between the two becomes a number product control reports rather than a
discovery.
\end{example}

\omcode{code/firm/platmap/firm_platmap.py}{85}{107}{Structural checks: one system of record per dataset, known names, one
producer per dataset, and no cycle.}\label{lst:pl:the-platform-map:validate}

\begin{dated}{2026-09}{How far banks are from one owner per number}\label{dat:pl:the-platform-map:bcbs}
The Basel Committee's principles for effective risk data aggregation and risk reporting (BCBS~239, published on 9 January 2013)
ask banks to generate accurate, complete and timely risk data from a data architecture that supports it. In its progress
report of 28 November 2023 (consulted September 2026), the Committee found that of the 31 global systemically important banks
it assessed, only two were fully compliant with all the principles.
\end{dated}

\section{Batch and stream}

\begin{definition}[Batch processing, stream processing, end-of-day batch]\label{def:pl:the-platform-map:batch}
\emph{Batch processing}\index{batch processing} runs a job on a complete, bounded dataset (a day of fills, a vendor file) and
produces its output once. \emph{Stream processing}\index{stream processing} runs a job continuously on an unbounded sequence
of events, updating its output as each event arrives. The \emph{end-of-day batch}\index{end-of-day batch} is the set of batch
jobs that run after the close to produce the day's official positions, marks, P\&L, risk and reports.
\end{definition}

A trading firm runs both, often for the same quantity. The position service of chapter~17 keeps positions as a stream during
the day, from each fill as it arrives, because traders and the risk gate need them now; the end-of-day positions job recomputes
them as a batch from the complete day, the clearing statement and the security master, because the books are closed on those.
The two must agree, and when they do not the difference is a break (chapter~22).

In a batch chain the questions are about time: when will each output be ready, which job decides when the chain ends, and how
late may each input be. With unlimited machines (chapter~13 removes that assumption), two passes over the map answer them.

\begin{method}[Forward and backward passes]\label{met:pl:the-platform-map:passes}
Order the jobs so that every job comes after the producers of its inputs.
\emph{Forward:} for each job, its start is the latest ready time of its inputs (or its fixed start), its end is start plus
duration, and its outputs are ready at its end; external inputs are ready at their arrival times.
\emph{Backward:} for each job in reverse order, its latest end is the earliest of its own deadline and the latest starts of the
jobs that read its outputs; its latest start is that minus its duration. The \emph{slack} of a job is its latest start minus its
scheduled start. The latest safe arrival of an external input is the earliest latest start among the jobs that read it.
\end{method}

\begin{proposition}[The dependency-driven schedule is the earliest]\label{prop:pl:the-platform-map:earliest}
With unlimited capacity and fixed durations, the forward pass of a map whose jobs all start on their inputs gives each job the
earliest start any schedule can give it, and the chain ends at the end of its longest chain of jobs: the chain that ends last,
traced back through the input that was ready last at each step.
\end{proposition}

\begin{proof}
By induction in dependency order. A job cannot start before all its inputs are ready; if each input is ready as early as
possible (the induction hypothesis, or its arrival time for an external input), the job starts as early as possible by
starting at the latest of those times. Following the last-ready input back from the job that ends last gives a chain of jobs
with no gap between one's end and the next's start, whose total is the end time.
\end{proof}

\Cref{lst:pl:the-platform-map:passes} is both passes. The forward pass also records every read of an input that was not yet
complete when the job started, which is only possible for a job with a fixed start: that is the check the hook needed.

\omcode{code/firm/platmap/firm_platmap.py}{185}{219}{The forward pass (start on the last input, or at the fixed time), the stale
reads it implies, and the backward pass of latest starts.}\label{lst:pl:the-platform-map:passes}

\begin{omfigure}
\begin{tikzpicture}
\begin{axis}[width=12cm,height=8.2cm,xbar stacked,bar width=6pt,xmin=0,xmax=16,
  xtick={0,2,4,6,8,10,12,14,16},xticklabels={16:00,18:00,20:00,22:00,00:00,02:00,04:00,06:00,08:00},
  ytick={0,1,2,3,4,5,6,7,8,9,10,11,12,13,14},yticklabels from table={figdata/platforms/01-the-platform-map/gantt_nominal.csv}{job},
  tick label style={font=\scriptsize},ymin=-0.7,ymax=14.7,grid=major,grid style={gray!20},
  legend columns=2,legend style={font=\footnotesize,at={(0.5,-0.12)},anchor=north,draw=none,/tikz/every even column/.append style={column sep=8pt}},
  table/col sep=comma]
\addplot[draw=none,fill=none,forget plot] table[x=start,y=k]{figdata/platforms/01-the-platform-map/gantt_nominal.csv};
\addplot[fill=omDef!45,draw=omDef] table[x=other,y=k]{figdata/platforms/01-the-platform-map/gantt_nominal.csv};
\addplot[fill=omThm!55,draw=omThm] table[x=chain,y=k]{figdata/platforms/01-the-platform-map/gantt_nominal.csv};
\draw[omThm,thick,dashed] (axis cs:15,-0.7) -- (axis cs:15,14.7) node[pos=0.97,left,font=\scriptsize]{07:00 deadline};
\legend{job,longest chain}
\end{axis}
\end{tikzpicture}

\omcaption{The miniature firm's \omterm{def:pl:the-platform-map:batch}{end-of-day batch} on a nominal night, every job started on its inputs. The longest chain (red)
runs from the corporate-action file through the golden copy of the reference data, the research features and the overnight
backtests, and ends at 03:55; the dashed line is the 07:00 deadline. Data: \texttt{fig\_platmap.py}.}\label{fig:pl:the-platform-map:gantt}
\end{omfigure}

On a nominal night (\cref{fig:pl:the-platform-map:gantt}) the chain ends at 03:55, 185 minutes before its deadline. Its longest
chain is not the risk batch, which is the longest job (six hours, from 20:15 to 02:15), but the research branch: the golden copy
waits for the 19:30 corporate-action file, the features wait for the golden copy, and six hours of backtests follow. A firm that
bought more machines for the risk grid would not move the end of the night by a minute. \Cref{tab:pl:the-platform-map:latest}
turns the backward pass into the question operations actually ask the vendor: how late can your file be?

\begin{table}[ht]
\centering\small
\begin{tabular}{lrrr}
\toprule
input & arrives & latest safe arrival & margin (min) \\
\midrule
fills, capture & 16:05 & 00:40, 22:25 & 515, 380 \\
closing prices & 16:20 & 00:40 & 500 \\
curve and volatility snapshot & 16:45 & 00:20 & 455 \\
vendor reference data & 18:00 & 22:35 & 275 \\
vendor corporate actions & 19:30 & 22:35 & 185 \\
clearing statement & 21:30 & 05:50 & 500 \\
\bottomrule
\end{tabular}
\caption{Arrival and latest safe arrival of each external input for every morning report to be ready by 07:00, from the
backward pass on the nominal map. The corporate-action file has the smallest margin. Data: \texttt{pl\_platmap.latest\_arrivals}.}\label{tab:pl:the-platform-map:latest}
\end{table}

\section{Reading a platform map}

Two more questions complete the reading. The first is structural: what does an input reach?

\begin{definition}[Blast radius]\label{def:pl:the-platform-map:blast}
The \emph{blast radius}\index{blast radius} of a dataset is the set of datasets derived from it, directly or through other
jobs: everything that is wrong if it is wrong and late if it is late.
\end{definition}

On the miniature firm's map the day's fills, the vendor
reference file and the corporate-action file each reach nine of the fifteen derived datasets; the clearing statement reaches
two (the reconciliation breaks and the signed P\&L); the market-data capture reaches four. A wrong corporate action is as
dangerous as a wrong fill, and much less watched: fills are reconciled with the venue's drop copy (Book~13, chapter~21) within
seconds, while a corporate action is checked, if at all, when a report looks odd.

The second question is about time, and it is the hook's. A job that starts on a fixed clock time reads whatever version of its
inputs exists at that moment. If an input's producer runs late, the job reads the previous version without an error, and its
output silently disagrees with outputs computed from the new one.

\begin{proposition}[When a fixed start reads stale data]\label{prop:pl:the-platform-map:stale}
A job with fixed start $s$ reads a complete version of input $x$ if and only if $x$ is ready by $s$. If $x$ is the output of a
chain of dependency-driven jobs of total duration $D_x$ that starts on an external input arriving at time $a$, the job reads
stale data exactly when $a > s - D_x$.
\end{proposition}

\begin{proof}
The first part is the definition of the forward pass. For the second, the forward pass gives $x$ the ready time $a + D_x$ when
$a$ is the last input of the chain, and the read is stale exactly when $a + D_x > s$.
\end{proof}

The risk batch of the hook starts at 21:00 ($s = 300$ minutes after the close). Its inputs include the positions, ready 45
minutes after the corporate-action file arrives (25 minutes of golden copy, 20 of positions), so it reads stale positions as soon
as the file arrives after 20:15. On a nominal night the file comes at 19:30, 45 minutes of margin, and nobody notices the fixed
start. On the night of the hook it came at 21:30: the golden copy was ready at 21:55, the positions at 22:15, and the risk batch
had been running on the old versions since 21:00 (\cref{fig:pl:the-platform-map:hook}). The P\&L job, dependency-driven, read
the new ones at 22:15. The three teams' morning was the price of one fixed start.

\begin{omfigure}
\begin{tikzpicture}
\begin{axis}[width=12cm,height=4.4cm,xbar stacked,bar width=8pt,xmin=4,xmax=12,
  xtick={4,5,6,7,8,9,10,11,12},xticklabels={20:00,21:00,22:00,23:00,00:00,01:00,02:00,03:00,04:00},
  ytick={0,1,2,3},yticklabels from table={figdata/platforms/01-the-platform-map/gantt_hook.csv}{job},
  tick label style={font=\scriptsize},ymin=-0.6,ymax=3.6,grid=major,grid style={gray!20},table/col sep=comma]
\addplot[draw=none,fill=none,forget plot] table[x=start,y=k]{figdata/platforms/01-the-platform-map/gantt_hook.csv};
\addplot[fill=omDef!45,draw=omDef] table[x=dur,y=k]{figdata/platforms/01-the-platform-map/gantt_hook.csv};
\draw[omThm,thick,dashed] (axis cs:5.9167,-0.6) -- (axis cs:5.9167,3.6);
\draw[omThm,thick,dashed] (axis cs:6.25,-0.6) -- (axis cs:6.25,3.6);
\node[font=\scriptsize,omThm,anchor=west] at (axis cs:6.35,1.55) {security master ready 21:55, positions 22:15};
\end{axis}
\end{tikzpicture}

\omcaption{The night of the hook: the corporate-action file arrives at 21:30. The risk batch, started by the clock at 21:00, has
already read the security master and the positions (their new versions are ready at the dashed lines); the P\&L job, started on
its inputs, reads the new ones. Data: \texttt{fig\_platmap.py} (\texttt{firm\_platmap.stale\_reads}).}\label{fig:pl:the-platform-map:hook}
\end{omfigure}

The fix is not to move the fixed start to 22:00, which only moves the threshold; it is to make the risk batch start on its
inputs. Dependency-driven, on the night of the hook it starts at 22:15 and ends at 04:15, well inside the deadline, and the map
reports no stale read. Fixed starts exist for good reasons (a machine is free only after 21:00; a vendor file has no completion
signal), and where one is kept, the map's check turns the silent disagreement into a failed precondition: the job refuses to
start on an incomplete input, or starts and marks its output as built on a stale version.

\begin{remark}[What the map does not know]\label{rem:pl:the-platform-map:limits}
The passes assume fixed durations and unlimited machines. Durations vary from night to night, jobs share a cluster, and a job
can fail and be rerun; chapter~13 adds queues and failures, and chapter~20 sizes the risk batch itself. The map also knows only
the flows someone declared: a spreadsheet that copies the marks at 18:00 is invisible to it, which is why the map is enforced by
making systems read their inputs through it (chapters~6, 7 and~16) rather than drawn beside them.
\end{remark}

\section{Tutorial: the miniature firm as a map}\label{tut:pl:the-platform-map:tut}

\begin{tutorial}
\textbf{Goal.} Describe the miniature firm as data, check it, schedule a night, and find the hook's defect by program.
\textbf{End state:} \cref{fig:pl:the-platform-map:gantt,fig:pl:the-platform-map:hook}, \cref{tab:pl:the-platform-map:latest}
and the \omterm{def:pl:the-platform-map:blast}{blast radius} of every input.
\begin{tutsteps}
\item \textbf{Declare.} Read \texttt{pl\_platmap}: systems by layer, datasets with their \omterm{def:pl:the-platform-map:sor}{system of record} and, for external
inputs, their arrival times, and the jobs of \cref{lst:pl:the-platform-map:jobs}.
\item \textbf{Validate.} \texttt{draft().validate()} reports that the marks have two systems of record; \texttt{firm().validate()}
reports nothing.
\item \textbf{Schedule.} \texttt{firm().schedule()} and \texttt{longest\_chain()} give \cref{fig:pl:the-platform-map:gantt}:
the chain golden copy, features, backtests, ending at 03:55. \texttt{slack(900)} and \texttt{latest\_arrivals} give
\cref{tab:pl:the-platform-map:latest}.
\item \textbf{Replay the hook.} \texttt{firm(risk\_start=300).stale\_reads(HOOK\_NIGHT)} returns the two stale reads of the risk
batch; \texttt{firm().stale\_reads(HOOK\_NIGHT)} returns none.
\item \textbf{\omterm{def:pl:the-platform-map:blast}{Blast radius}.} \texttt{firm().blast\_radius()} and \texttt{downstream} of the corporate-action file.
\end{tutsteps}
\textbf{What to change next.} Give the backtests a deadline of 06:00 and find the new latest arrival of the corporate-action file;
add a margin-forecast job that reads the positions and the clearing statement and must be ready by 06:00, and see which input's
latest arrival it moves.
\end{tutorial}

\section{Build: the platform map}\label{bld:pl:the-platform-map:platmap}

\begin{build}
\bfield{Purpose} The firm's map of systems, datasets and jobs as data, with the checks and the time questions every later
chapter uses: the tick store (chapter~4) and the reference data (chapter~6) declare their outputs in it, the data-quality rules
of chapter~7 and the observability of chapter~28 attach to its datasets, and the scheduler of chapter~13 runs its jobs.
\bfield{Interface} \texttt{System(name, layer, owner)}, \texttt{Dataset(name, system, cadence, ready)}, \texttt{Job(name,
inputs, outputs, duration, start, deadline)}; \texttt{PlatformMap(systems, datasets, jobs)} with \texttt{validate},
\texttt{producer}, \texttt{downstream}, \texttt{upstream}, \texttt{schedule}, \texttt{ready\_times}, \texttt{stale\_reads},
\texttt{latest\_starts}, \texttt{slack}, \texttt{longest\_chain}, \texttt{blast\_radius}; times in minutes after the close.
\bfield{Rules} One \omterm{def:pl:the-platform-map:sor}{system of record} per dataset and one producer per derived dataset; no cycles; a fixed-start job is allowed
but every read before its input is ready is reported; the map is the only place a flow is declared.
\bfield{Acceptance tests} \texttt{code/firm/platmap/tests/}: the forward pass and the longest chain on a small map, and how a
late input moves both; the backward pass and slack by hand; a stale read found and absent; two systems of record, an unknown
system, an orphan dataset and a cycle reported; closures and \omterm{def:pl:the-platform-map:blast}{blast radius}.
\bfield{Stretch} Durations as distributions and the probability of meeting the deadline by simulation; versions of datasets
and the version each job read; a map generated from the systems' own declarations rather than written by hand.
\end{build}

\begin{omsources}
\item Basel Committee on Banking Supervision, \emph{Principles for effective risk data aggregation and risk reporting}
(BCBS~239), January 2013, and \emph{Progress in adopting the Principles}, November 2023.
\item T.~Akidau et al., ``The dataflow model'', \emph{Proceedings of the VLDB Endowment} 8(12), 2015 (batch as a special case
of stream processing).
\end{omsources}

\section{Exercises}

\begin{exercise}[$\star$]\label{exo:pl:the-platform-map:1}
The P\&L job reads the positions (ready at 20:15), the marks (17:05) and the security master (19:55), and takes 25 minutes.
When does it start and end on a nominal night, and which input decides?
\end{exercise}

\begin{exercise}[$\star$]\label{exo:pl:the-platform-map:2}
The vendor's reference-data file arrives at 22:50 instead of 18:00. Using \cref{tab:pl:the-platform-map:latest}, which report is
late, and by how many minutes?
\end{exercise}

\begin{exercise}[$\star$]\label{exo:pl:the-platform-map:3}
Why does the clearing statement have the smallest \omterm{def:pl:the-platform-map:blast}{blast radius} of the seven inputs, and does a small \omterm{def:pl:the-platform-map:blast}{blast radius} mean the
input needs less checking?
\end{exercise}

\begin{exercise}[$\star\star$]\label{exo:pl:the-platform-map:4}
To finish earlier, the risk team moves the risk batch's fixed start from 21:00 to 20:30. Up to what arrival time of the
corporate-action file does it read complete inputs, and what margin does that leave on a nominal night?
\end{exercise}

\begin{exercise}[$\star\star$]\label{exo:pl:the-platform-map:5}
Product control's independently verified marks differ from the pricing library's by a few basis points on some positions.
Explain why these are two datasets rather than two copies of one, and which of the two the P\&L sign-off should read.
\end{exercise}

\begin{exercise}[$\star\star$]\label{exo:pl:the-platform-map:6}
During the day the position service keeps positions as a stream; after the close a batch job recomputes them. Give two reasons
the two can differ at 17:00, and say which one the books are closed on.
\end{exercise}

\begin{exercise}[$\star\star\star$]\label{exo:pl:the-platform-map:7}
\emph{Coding.} Add a job \texttt{margin forecast} that reads the positions and the clearing statement, takes 60 minutes and must
be ready by 06:00. What is its latest start, does it change any other job's latest start, and what is the clearing statement's
new latest safe arrival?
\end{exercise}

\begin{exercise}[$\star\star\star$]\label{exo:pl:the-platform-map:8}
\emph{Find the flaw.} ``The overnight batch is too close to its deadline. We will double the risk grid's machines, which halves
the risk batch to three hours, and the whole batch will end three hours earlier.''
\end{exercise}

\section{Problem: The Morning the Numbers Disagreed}

\begin{problem}[{Weekend problem --- a late file and a fixed start}]\label{pb:pl:the-platform-map:1}
The miniature firm's map (\texttt{pl\_platmap}), its nominal night and the night the corporate-action file arrived at 21:30.

\medskip\noindent\textbf{Part I --- The nominal night.}
\begin{enumerate}
\item When do the golden copy, the positions and the risk batch start and end, with every job started on its inputs?
\item Which job ends last, and when?
\item What is the longest chain, and why is the risk batch, the longest job, not on it?
\item How many minutes of slack does the batch have before 07:00?
\item What is the risk batch's slack?
\end{enumerate}

\medskip\noindent\textbf{Part II --- How late can the files be?}
\begin{enumerate}[resume]
\item What is the latest safe arrival of the vendor's corporate-action file, and which job sets it?
\item What is the latest safe arrival of the market-data capture, and why is it earlier than that of the fills?
\item Which input has the largest margin between its nominal and its latest safe arrival?
\item How many datasets does the corporate-action file reach?
\item Which datasets does a wrong clearing statement reach?
\end{enumerate}

\medskip\noindent\textbf{Part III --- The night of the hook.}
\begin{enumerate}[resume]
\item The file arrives at 21:30. When are the security master and the positions ready?
\item The risk batch starts at 21:00 by the clock. Which of its inputs does it read stale?
\item When does the P\&L job run, and which versions does it read?
\item Up to what arrival time of the file would the fixed-start risk batch have read complete inputs?
\item Started on its inputs, when would the risk batch have run on that night, and would it have met the deadline?
\end{enumerate}

\medskip\noindent\textbf{Part IV --- The verdict.}
\begin{enumerate}[resume]
\item State the \emph{named result}: the end and slack of the nominal night, the latest safe arrival of the corporate-action
file, and the arrival time beyond which the fixed-start risk batch reads stale data.
\item Why did the disagreement produce no error in any system?
\item Which of the three failure modes of \cref{prop:pl:the-platform-map:oneowner} occurred?
\item Give two ways to keep a fixed start without the silent disagreement.
\item In one sentence: what does a platform map give that a system diagram does not?
\end{enumerate}
\end{problem}

\section{Interview questions}

\begin{interviewq}[$\star$ \iqroles{developer, risk}]\label{iq:pl:the-platform-map:1}
What is a \omterm{def:pl:the-platform-map:sor}{system of record}, and why does a firm need one per dataset?
\end{interviewq}

\begin{interviewq}[$\star$ \iqroles{developer}]\label{iq:pl:the-platform-map:2}
When would you compute positions as a stream and when as a batch?
\end{interviewq}

\begin{interviewq}[$\star\star$ \iqroles{developer, risk}]\label{iq:pl:the-platform-map:3}
The morning risk report and the P\&L disagree on one position. How do you find out why?
\end{interviewq}

\begin{interviewq}[$\star\star$ \iqroles{developer}]\label{iq:pl:the-platform-map:4}
Should a batch job start at a fixed time or when its inputs are ready? What goes wrong with each?
\end{interviewq}

\begin{interviewq}[$\star\star$ \iqroles{developer}]\label{iq:pl:the-platform-map:5}
How would you find out how late a vendor file can arrive before the morning reports are late?
\end{interviewq}

\begin{interviewq}[$\star\star\star$ \iqroles{developer, risk}]\label{iq:pl:the-platform-map:6}
You join a firm with no diagram of its systems. How do you build a map that stays correct, and what do you check first?
\end{interviewq}
